\chapter{Conclusion and Future Work}
\label{chp:conclusion}

% \vfill

% % This thesis set out to establish what \gls{spar} costs to run.

% % The first section below states what the measurements show about the protocol's practicality and where its limits lie. The second collects the optimizations and extensions that the implementation leaves open, several of which we expect to change the picture materially.
% \newpage

\section{Conclusion}
This thesis set out to answer two questions about \gls{spar}. What does the protocol cost to run, and how far does it scale before it becomes impractical.

On the first question, the measured cost is substantially higher than the original paper anticipates. The cost of a round is dominated by the server's write step, and the write step is dominated by internal products, an operation absent from the cost analysis in~\cite{cryptoeprint:2025/860}. A single internal product costs roughly twenty times an external product at our parameters (\autoref{tab:bench:rgsw}). The paper's expectation of round times on the order of seconds for a hundred users is not feasible with this instantiation.

On the second question, computation time, not noise, is the binding constraint. A single \texttt{server::Write} call grows linearly in $\eta$ at roughly $4.2$ seconds per participant on our hardware, and since the server performs one call per client, the total server time per round grows quadratically. At $\eta = 16$ the server already spends around $18$ minutes per round, and extrapolating to $\eta = 100$ gives round times on the order of hours. The noise budget would allow the protocol to run somewhat further, with an estimated $\eta < 32$ at $\ell = 3$ under multi-party noise flooding (\autoref{tab:noise:multiparty}), and the protocol itself contributes so little noise that a single-key instantiation supports a practically unbounded number of users. Time renders the protocol impractical well before either limit matters.

The positive result is that \gls{spar} can be instantiated at all from well-established \gls{fhe} primitives. The protocol runs end-to-end on standard \gls{rlwe} and \gls{rgsw} constructions in OpenFHE, with correct decryption under multi-party thresholdized keys. The obstacles we encountered were quantitative rather than structural, which suggests the optimizations below could change the picture materially.

\section{Future Work}
The most direct algorithmic lever is \gls{simd} packing (\autoref{def:simd}) in the write loop~\cite{cryptoeprint:2011/133}. Packing lets the server evaluate $N$ rounds simultaneously, dividing the per-round cost of the dominant loop by $N$, though converting this into throughput requires protocol changes on the client side as discussed in \autoref{sec:results:discussion}.

On the implementation side, OpenFHE's AVX512 backend was left unmeasured for lack of suitable hardware (\autoref{sec:avx512}). Quantifying its concrete speedup, and choosing moduli sized to fill the $512$-bit registers, is a natural next step.

% A bandwidth optimization using \gls{rgsw} or packed \gls{rlwe} ciphertexts~\cite{chillotti:onionring:2019, cryptoeprint:2025/1088} was originally planned but dropped due to time constraints. It would reduce what each client uploads, at the price of a more involved setup phase since the server requires $\rgsw{s}$.

Finally, our implementation covers the core of \gls{spar} but not the full protocol, which additionally specifies equivocation protection, bandwidth optimizations, and the sorting of messages into buckets. Completing these components is required before the end-to-end system can be evaluated, but would do little to improve the cost of the protocol.
